\section{Introducción}

\subsection{Contexto del trabajo}

La Región de La Araucanía concentra actividad agrícola, agroindustrial y de pequeñas y medianas
empresas productivas, con instalaciones frecuentemente dispersas: predios, invernaderos, bodegas,
cámaras de frío y estanques distribuidos en el territorio. Estas condiciones hacen del monitoreo
remoto una capacidad de valor operativo directo, y de las tecnologías de radio de largo alcance y bajo
consumo una respuesta técnicamente apropiada.

En los últimos años, la adquisición de sensores dejó de constituir una barrera económica relevante. Un
nodo sensor con radio de largo alcance tiene un costo de entre 5 y 30 dólares, y la banda de
915--928\,MHz que Chile emplea bajo el plan regional AU915 es de uso libre y no requiere licencia de
espectro. El resultado es que la dificultad se desplazó: ya no está en obtener el sensor, sino en qué
hacer con lo que el sensor emite.

\subsection{Propósito de este documento}

Este informe corresponde a la \textbf{\datoEtapa} del \datoCurso{} \datoCodigo, cuyo alcance definido
por la pauta de evaluación comprende la definición del problema, la formulación de objetivos, el
análisis de requerimientos y la planificación mediante un backlog priorizado. En términos de la
metodología adoptada, corresponde a las fases \emph{Prepare} y \emph{Plan} del ciclo PPDIOO.

\subsection{Estructura del informe}

La sección~\ref{sec:problema} define el problema mediante un árbol de causas y efectos, justifica su
relevancia y delimita el alcance del trabajo. La sección~\ref{sec:objetivos} formula el objetivo
general y los objetivos específicos con sus criterios SMART y la matriz de trazabilidad que los vincula
al problema. La sección~\ref{sec:metodologia} expone y justifica el marco metodológico. La
sección~\ref{sec:requerimientos} presenta el catálogo de requerimientos y la evaluación crítica de
alternativas de solución. La sección~\ref{sec:protocolo} desarrolla en detalle la decisión de mayor
peso técnico del proyecto: el protocolo de comunicación. La sección~\ref{sec:planificacion} presenta el
backlog y la verificación de viabilidad del alcance. La sección~\ref{sec:conclusiones} recoge las
conclusiones de la etapa y el trabajo pendiente.

\begin{cajaaviso}[title={Estado del documento}]
Este informe se entrega como avance de etapa y contiene dos elementos declaradamente incompletos, ambos
señalados en el texto donde corresponde:

\begin{enumerate}
\item \textbf{El levantamiento directo con las partes interesadas está pendiente.} Las declaraciones de
la organización destinataria se incorporarán a la sección~\ref{sec:interesados} y al catálogo de
requerimientos una vez realizadas las entrevistas. Hasta entonces, los requerimientos atribuidos al
propietario del proyecto se sustentan en las indicaciones transmitidas por el profesor guía.

\item \textbf{Las fechas de las cuatro etapas están sujetas a confirmación}, conforme se explica en la
sección~\ref{sec:calendario}.
\end{enumerate}

Ambos elementos se declaran explícitamente en lugar de omitirse, y su resolución constituye la acción
prioritaria del período en curso.
\end{cajaaviso}
